% In this file you should put the actual content of the blueprint.
% It will be used both by the web and the print version.
% It should *not* include the \begin{document}
%
% If you want to split the blueprint content into several files then
% the current file can be a simple sequence of \input. Otherwise It
% can start with a \section or \chapter for instance.

\chapter{Basic Examples}

We start with a very simple definition.

\begin{definition}\label{def:my_square}
\lean{my_square}
\leanok
Let $x$ be a real number. We define the square of $x$ as $x \times x$.
\end{definition}

Notice that we placed the \texttt{\textbackslash leanok} macro above, which will color this definition green in the graph because the Lean code is complete.

\begin{lemma}\label{lem:square_nonneg}
\uses{def:my_square}
\lean{square_nonneg}
\leanok
For any real number $x$, we have $0 \le x^2$.
\end{lemma}
\begin{proof}
\leanok
This follows trivially from the properties of real numbers.
\end{proof}

Because we included \texttt{\textbackslash uses\{def:my\_square\}}, the blueprint graph will draw a dependency arrow from our definition to this lemma.

\begin{theorem}\label{thm:fermat}
\lean{fermat_last}
For any integers $x, y, z$ and $n > 2$, $x^n + y^n \ne z^n$.
\end{theorem}
\begin{proof}
The proof is too large to fit in this margin.
\end{proof}

Notice we did \textbf{not} include \texttt{\textbackslash leanok} for the theorem or its proof, because we used a \texttt{sorry} in our Lean file! This will show up as a blue (unverified) node in your graph.